\section*{Change Log}
\phantomsection
\addcontentsline{toc}{section}{Change Log}
\label{sec:change-log}

This page records the analysis-level changes between circulated versions of the note.
It is not a substitute for the configuration and figure provenance given in the body
and appendices. Numerical results from different versions should not be combined: each
version corresponds to a distinct, internally consistent analysis state.

\begingroup
\footnotesize
\setlength{\tabcolsep}{4pt}
\renewcommand{\arraystretch}{1.04}
\setlength{\LTleft}{\fill}
\setlength{\LTright}{\fill}
\begin{longtable}{P{0.11\textwidth}P{0.16\textwidth}P{0.64\textwidth}}
\caption{Analysis-note version history. The v4.2 row defines the accepted
analysis result and the subsequent rows define the additional diagnostic scope
of the present document. Earlier versions are summarized to keep changes in
inputs, validation, and statistical interpretation traceable to reviewers.}
\label{tab:note-change-log}\\
\toprule
Version & Circulation date & Principal changes \\
\midrule
\endfirsthead
\multicolumn{3}{l}{\footnotesize\itshape Table~\ref{tab:note-change-log}
continued from the previous page.}\\
\toprule
Version & Circulation date & Principal changes \\
\midrule
\endhead
\midrule
\multicolumn{3}{r}{\footnotesize\itshape Continued on the next page.}\\
\endfoot
\bottomrule
\endlastfoot
v1 & March--May 2026 &
Established the full-range Gaussian-process background model, the blinded local
signal-extraction likelihood, the radiative-fraction conversion to $\eps^2$, and the
first single-dataset and combined validation studies. The figures in this version were
development results and did not define the present candidate configuration. \\
\addlinespace
v2 & June 2026 &
Introduced the $2.25\,\sigma_m$ blind/training geometry, matched-reference full-refit
pseudoexperiments, the shared-$\eps^2$ simultaneous likelihood, and expanded 2015,
2016, and 2021 method documentation. This version used common or development-stage
length-scale choices and retained broader engineering-run scan ranges in several
studies. \\
\addlinespace
v3 & July 10, 2026 &
Uses the regenerated 2021 1\% high-$p_{\mathrm{sum}}$ histogram, with no vertex-fit
$\chi^2$ requirement and at least eight electron-track hits; restricts the 2015 and
2016 scans to \SIrange{19}{90}{MeV} and \SIrange{39}{180}{MeV}; freezes the
dataset-dependent lower length-scale factors at 1.0, 0.9, and 1.1; and adopts 90\% CL
upper limits obtained with \CLs{}. The main validation section now contains the
functional-form seed construction and the injected-signal lower-bound scans, while
superseded closure and limit studies are retained with their original context in the
appendices. The results section documents the individual and simultaneous limits,
local and approximate global $p$-values, and all remaining pre-unblinding audit
qualifications. \\
\addlinespace
v4 & August 3, 2026 &
Separates the GP fit support from the signal-search interval in all three campaigns.
The 2015 and 2016 searches keep their reviewed mass grids but regain sidebands beyond
both endpoints, while the 2021 10\% sample retains the wider support developed in the
recent endpoint study. The result now uses the full 2015 and 2016 datasets in the
shared-$\eps^2$ fit and includes conditional background-only expected bands. The same
pseudoexperiments place the observed limit within the ensemble through the strong-,
weak-, and bounded two-sided consistency tests used in v3. Those quantities remain
upper-limit diagnostics rather than discovery significances or a substitute for
full-procedure coverage. \\
\addlinespace
v4.1 & August 4--5, 2026 &
Responds to the full-2016 length-scale saturation with a controlled observed-only
scan of $k_{\max,2016}=8,10,12,15,$ and 20. Factor 12 is the first nonbinding value
and is followed by a stable plateau. The exact combined observed limit and local
asymptotic $p_0$ scan are recomputed with all other analysis settings fixed. A
ten-toy 2021 source-and-exposure screen---native 1\%, $1\%\times10$,
$1\%\times100$, native 10\%, and $10\%\times10$---is included only as a fit-space
hyperparameter pilot. Its reviewed background-only scan supports factor 20 for the two
projected-100\% lanes and factor 25 if one bound must cover every pilot exposure;
the fixed-amplitude screen shows no factor-20 sensitivity loss, while a common 2--3\%
under-response remains an unresolved ten-toy closure diagnostic. No new expected
bands, limit-tail ensemble result, toy-calibrated discovery claim, or coverage claim
is included; the observed result remains a candidate pending production-matched
closure and direct coverage. The August 5 methodology revision adds explanatory
schematics for $C$, $\ell$, fixed $\alpha_i$, and the per-hypothesis optimization;
it changes no configuration or numerical result. \\
\addlinespace
v4.2 & August 5, 2026 &
Accepts $k_{\max,2016}=12$ as the current full-2016 setting and completes the
232-mass combined 90\% \CLs{} search with exactly 300 conditional background-only
joint pseudoexperiments at each mass. Central 68\% and 95\% intervals are
reported for that principal shared-$\eps^2$ search. The note also adds
standalone observed limits for the individual 2015 full, 2016 full, and 2021
10\% distributions, auxiliary 100-pseudoexperiment standalone and pairwise
conditional limit bands, the combined strong-, weak-, and bounded two-sided
limit-tail diagnostics, and combined and constituent local asymptotic $p_0$
plots. The auxiliary ensembles select indices 0--99 from the accepted
300-draw parent stream and reuse campaign pseudo-spectra across scopes. No
scan-toy-calibrated global discovery probability or direct coverage claim is
introduced; all bands remain descriptive quantiles conditional on the reviewed
fixed GP states. The results section also
adds an exact fixed-state 65~MeV extraction, including sideband context,
background-subtracted spectra, conditional Pearson-residual diagnostics, and a
comparison of the three standalone signal estimates with the simultaneous
shared-$\eps^2$ fit. \\
\addlinespace
v4.2 follow-up & August 6, 2026 &
Adds the observed-equivalent comparison with the published BaBar visible
90\% limit while scaling only the 2021 count-density response from 10\% to
100\% statistics.  It adds an exact fixed-mass 0.5~MeV common-binning refit
and 1.25~MeV stress test after documenting why exact 0.625~MeV aggregation is
not possible from the histogram-only 2015 and 2016 inputs.  The Figure~62
coefficient panel now uses the physical-domain nominal asymptotic profile set
instead of a symmetric signed-Wald bar.  A separate appendix records two
fixed-seed, background-only conditional replacements of the 2021 central
window, their complete observed/asymptotic scans, and optimizer-repeat
ledgers.  A deterministic central-mean follow-up separates the functional
truth/GP-analysis shoulder below 65~MeV from the additional fluctuation in its
fixed draw.  Ten paired fixed-GP-mean replacements compare the exactly
equivalent 2.25/2.5-$\sigma_m$ production-bin scope with a 3-$\sigma_m$
scope, reporting only descriptive observed-limit and local-$p_0$ summaries.
This addendum changes no accepted v4.2 card or principal result, adds no
projected bands or projected $p$-values, and introduces no coverage or
scan-calibrated global-significance claim. \\
\addlinespace
v4.2 candidate catalogue & August 7, 2026 &
Defines an operational eight-feature catalogue from the three standalone,
three genuine-overlap pairwise, and all-three v4.2 upper-limit curves.  For
each feature and active dataset, complete analysis bins within the
dataset-specific $\pm2.5\sigma_m$ interval are replaced in the native integer
histogram by one deterministic GP-posterior-mean lane and three joint
GP-posterior plus Poisson conditional inpainting replicates.  The accepted
v4.2 moving-mask fit and shared-$\eps^2$ combination are rerun at nearby
candidate positions, including a dense reciprocal 65/80~MeV test.  The
appendix adds common-axis 65~MeV positional overlays and selected-row
kinematic distributions from the available v11 2021 10\% SIMP/TM atlas.  The
atlas is not the exact v16 prompt-TC sample used by the observed limit.  This
addition is a data-selected mechanism study; it changes no accepted observed
limit and makes no coverage, global-significance, or candidate-promotion
claim. \\
\addlinespace
v4.3 exact-v16 diagnostics & August 10, 2026 &
Adds the exact-v16 2021 10\% topology, stored-trigger, and mass-versus-control
histograms while leaving the accepted v4.2 mass-only search unchanged.  The
three canonical diagnostic baselines reproduce the frozen v4.2 spectrum
bitwise, and all saved topology and trigger partitions close bin-by-bin.  A
zero-toy, histogram-level conditional-composition GPR masks the union of the
65 and 80~MeV stripes and selects its RBF card only through deterministic
fake-gap closure.  Mass-binning, control-grouping, and single-target-mask
variations are retained as model-stability checks.  The addendum reports no
event-level 2D inference, signal yield, significance, category limit,
expected band, or coverage result.  Joint masking of the 65 and 80~MeV
regions is a guard against cross-response in their shared GP/background
profile and adjacent deficit/background structure.  The response is assigned
to the common fitting procedure, with no physical-dependence or shared-event
interpretation. \\
\addlinespace
v4.3 region comparisons & August 11, 2026 &
Adds exact-v16 central-window, lower-sideband, and upper-sideband comparisons
for momentum sum, vertex-projection significance, minimum track $|y_0|$,
electron 3D-hit count, exhaustive topology, and stored-trigger category at
51, 65, 80, 90, 117, and 162~MeV.  Region totals use fractional overlap of
the 0.125-MeV mass-bin edges, and continuous-variable group boundaries are
frozen with the union of all six central regions masked.  The addition is
descriptive and histogram-level; it introduces no fitted signal yield,
significance, or limit.  The largest shape evolution occurs at the 51~MeV
lower boundary and the largest interior differences occur at 65~MeV; the
80~MeV differences are smaller. \\
\addlinespace
v4.5 refmatched exposure closure & August 11, 2026 &
Reuses only the reviewed Section~5.4 background pseudo-histograms for a new
frozen-v4.2, full-refit, matched-reference injection--extraction screen.  The
common \texttt{fGenGammaThresh} stress family is used for $1\%\times10$,
native 10\%, $1\%\times100$, and $10\%\times10$ at 65, 90, 120, 180, and
210~MeV with 0-, 1-, 3-, and 5-$\sigma_A$ injections and ten toys per cell.
The addition reports pull moments, signed $\eps^2$ coordinates, direct unpaired
source/exposure comparisons, factor-15 boundary contacts, and retained fit
outliers.  A separate deterministic product supplies toys 10--99 for later
continuation while reproducing the original ten bit-for-bit; those new toys do
not enter v4.5.  The study is a conditional screening diagnostic, not coverage,
an expected-limit ensemble, a production-card qualification, or a change to the
accepted v4.2 observed result. \\
\addlinespace
v4.6 full-100 refmatched closure & August 12, 2026 &
Uniformly reruns toy indices 0--99 for the four headline source/exposure
ensembles under the frozen v4.2 2021 card and a predeclared adaptive
optimizer-repeat gate.  Every attempt uses the identical cached pseudo-data,
an optimizer-only seed, 12 randomized restarts plus the nominal start, and
archived LML, hyperparameter, uncertainty, covariance, and warning diagnostics.
Branch selection uses only maximum GP LML and repeat agreement, never pull size
or sign.  Six irreproducible fit states are excluded from 8000; all 80 cells
retain 99 or 100 toys.  The independent toys-10--99 reserve confirms the
pilot-selected $1\%\times10$, 65~MeV zero-signal offset, while deterministic
analytic-mean closure identifies it as a local stress-truth/GP mismatch rather
than an optimizer or outlier effect.  The v4.1 and v4.5 ten-toy studies move to
an archival appendix.  The study remains conditional injection--extraction
closure, not direct \CLs{} coverage, an expected-limit ensemble, a global
significance calibration, or a change to the accepted v4.2 result. \\
\addlinespace
v4.7 near-threshold background modeling & August 12, 2026 &
Adds the section ``Validation of near-threshold background modeling.''
Independent logistic-turn-on times exponentiated-Chebyshev models are fitted
to the native 2021 1\% and 10\% spectra over 35--80~MeV; degree five is the
lowest common candidate passing the fixed goodness-of-fit gates.  Twenty
independent Poisson background spectra per source-derived scenario are tested
at 55, 60, 65, and 70~MeV with 0-, 1-, 3-, and 5-$\sigma_A$ injections under
the unchanged v4.2/v4.5 analysis semantics.  A pull-blind audit of the first
complete optimizer ledger exposed three additional displaced branches, so
the reference-relative repeat gate was amended before results were accepted
and all 40 task products were uniformly rerun.  The dedicated truth partially
reduces the $1\%\times10$ 65~MeV offset but does not give uniform closure:
native 10\% retains a positive 65~MeV mean, while both scenarios show negative
55~MeV means and $1\%\times10$ shows a positive 70~MeV mean.  The 20-toy
results are correlated conditional diagnostics, not direct coverage, expected
bands, observed limits, a causal comparison with the unpaired v4.6 ensemble,
or a change to the accepted v4.2 result. \\
\addlinespace
v4.8 residual-structured stress tests & August 14, 2026 &
Adds two post-exploratory 2021 generating means motivated by the v4.7 residual
structure: a fixed-node natural log spline selected between predeclared K2 and
K3 candidates by blocked source-only folds, and a fixed smoothly overlapping
low/middle/high functional blend.  Both forms fail the support-wide
source-qualification and signal-influence gates and are therefore retained
only as requested conditional stress truths.  Twenty nested Poisson
backgrounds per model and each of five explicit source/exposure lanes are
extracted at 65, 90, 120, 180, and 210~MeV with 0-, 1-, 3-, and
5-$\sigma_A$ injections, giving 4000 accepted states and no exclusions.  The
K2 branch has three exploratory material zero-signal flags versus sixteen for
the regional blend, but 345/400 K2 native-1\% states contact the pilot's
factor-25 length ceiling.  A pull-blind post-closure addendum with eight fresh
backgrounds qualifies factor 50 against a factor-75 sentinel for that targeted
lane only; the original closure is not rerun or uncensored, and no common
all-lane ceiling or production-card change is claimed.  The addition provides
no coverage, expected band, observed-data bias, limit, exclusion, or evidence
that either source family cannot absorb signal. \\
\addlinespace
v4.9 threshold-support qualification & August 17, 2026 &
Completes a paired one-factor diagnostic of a 30--300~MeV 2021 GP support while
retaining the 50--250~MeV search interval and v4.2/v4.5 extraction settings.  Exactly
25 full 30--300~MeV Poisson backgrounds in each of the $1\%\times10$ and native-10\%
source lanes are evaluated at 55, 60, 65, and 70~MeV with 0-, 1-, 3-, and
5-$\sigma_A$ injections under both 30--300 and 40--300~MeV GP supports.  All 1600
states are accepted with no exclusions or bound contacts.  Lowering the support has a
large paired effect at 65~MeV, but six of eight primary zero-signal cells are material
flags and the source-conditioned \texttt{fSigPowExpQ}-anchored truth is not globally
qualified.  Historical v4.6--v4.8 batteries move to an archival appendix and the
replacement-window ensemble is reduced to provenance.  No production freeze,
coverage, expected band, observed-data bias measurement, limit, or exclusion follows. \\
\addlinespace
v4.9.1 background-validation consolidation & August 17, 2026 &
Restores a comprehensive Section~5 while retaining the exhaustive v4.6--v4.9 study
batteries in appendices.  The historical four-lane, five-mass full-refit suite uses
the independently source-fitted \texttt{fGenGammaThresh} generator.  At 65~MeV only,
the historical $1\%\times10$ and native-10\% cells are replaced---not pooled---with
100-background threshold-aware ensembles: the Table~17 turn-on truth with 40--300~MeV
support and the \texttt{fSigPowExpQ}-anchored native-10\% truth with 30--300~MeV
support.  All 800 new injection states are accepted with zero exclusions.  Both
replacement zero-signal 90\% mean intervals lie inside the post-result
$|\overline p|<0.5$ practical band; all 20 composite zero-signal and all 80
lane--mass--strength point means also lie inside it.  Review-stable Figs.~64--66 show
the revised spurious-signal, targeted Table~17, and $2\times4$ mean/width summaries.
The note separately retains the targeted and composite Holm results, the three of
twenty zero-signal and fifteen of eighty all-strength unit-width incompatibilities,
stable 96.6--96.9\% injection response, and nonzero fixed-hyperparameter high-mass
tail influence.  Practical source-conditioned mean-bias validation is complete; no
truth-independent zero-bias proof, exact pull calibration, direct coverage,
observed-data bias measurement, limit, or exclusion is inferred. \\
\addlinespace
v4.9.5 2021 GP-support-edge optimization & August 20, 2026 &
Scans lower GP-support edges of 30, 32, 34, 36, 38, and 40~MeV against the same 100
archived native-10\% conditional backgrounds at 55, 60, 65, and 70~MeV with matched
0-, 2-, and 5-$\sigma_A$ injections.  Under the original rule, every candidate has
at least one injection cell with $|\overline p|\geq0.5$ whose 90\% interval excludes
zero, so no edge qualifies.  A documented post-phase-1 amendment
formalizes ``generally below 0.75'' and provisionally selects 36~MeV from the initial
25-background cohort; an independent 75-background continuation at 34, 36, and
38~MeV confirms 36--300~MeV without retuning.  A 201-point observed 2021 native-10\%
scan then reports aligned signal-yield limits, $\eps^2$ limits, and local analytic
$p_0$ values.  No upper-limit bands, \CLs{} toys, direct-coverage claim, or
scan-global significance calibration is introduced. \\
\addlinespace
v4.9.6 editorial revision & September 2, 2026 &
Revises the presentation of the validated v4.9.5 analysis without changing any data
input, analysis card, fit, figure, or numerical result.  The 2021 support-selection
study now closes Section~\ref{sec:toys}, and the corresponding observed scan now
closes Section~\ref{sec:results}.  Stale descriptions of the source samples and GP
support are corrected so that historical configurations remain distinct from the
selected native-10\% 36--300~MeV support.  The surrounding prose, section flow, and
typography are edited for clarity and consistency.  This is an editorial release,
not a new statistical analysis or result state. \\
\addlinespace
v4.9.7 full-2016 support qualification & September 2, 2026 &
Uses the pre-existing 2016 10\% development subset for partial observed-shape
information and the full-2016 spectrum only for one scalar 26--210~MeV normalization
before the support decision.  The hybrid generating mean contains an explicitly
waived broad-tail continuation whose archived record has
\texttt{fit\_ok=false}; it is used only as a conditional stress truth.  A frozen
25-background scan tests lower support edges from 28 through 34~MeV at four masses and
three injection strengths.  Every edge has only 3 of 12 cell means, including only
1 of 4 background-only means, below 0.75 in magnitude, and all violate the 1.25
gross-bias guard.  Two lanes also retain one technical exclusion each.  The declared
rule therefore returns no provisional edge and stops without retuning; no continuation,
full-2016 observed scan, updated combination, or 100-toy combined band is produced.
A separate post-unblinding audit shows that the earlier 65~MeV narrow-signal
interpretation is not robust to the 2021 GP-support prescription.  No direct coverage,
calibrated sensitivity, scan-global significance, or signal/background classification
is introduced. \\
\end{longtable}
\endgroup
